\documentclass[../main.tex]{subfiles}
\begin{document}
%==========================================TIÊU ĐỀ TẬP========================================
\begin{table}[h]
  \begin{adjustwidth}{-1cm}{}
    \begin{tabular}{>{\centering\arraybackslash}p{6cm}|>{\raggedright\arraybackslash}p{12.5cm}}
      \multirow{5}{6cm}
      % Cột bên trái: ÉPISODE và số tập
      {\\[-40pt]
        \flaregothic\fontsize{20pt}{20pt}\selectfont \centering
        \color{eptype}ÉPISODE\color{black}\\
        \burbank\fontsize{72pt}{72pt}\selectfont
        \color{epnum}104\color{black}\\[6pt]
        \cafeteria\fontsize{20pt}{20pt}\selectfont\color{epnum}(8-8)\color{black}
        \\[18pt]
        \flaregothic\fontsize{18pt}{18pt}\selectfont
        \color{epattr}LE PLUS POPULAIRE ! \\
        \color{black}
      }
       &
      % Dòng thứ nhất: tên tập tiếng Nhật
      {\fotskip\fontsize{18pt}{18pt}\selectfont \ruby{五}{ご}\ruby{月}{がつ}\ruby{病}{びょう}のあなたへ
      }                       \\[6pt]
      % Dòng thứ hai: tên tập tiếng Anh
       & {\cafeteria\fontsize{18pt}{18pt}\selectfont Ayame How-To: Beware, Beware the Blues of May}            \\[4pt]
      % Dòng thứ ba: tên tập tiếng Việt
       & {\vnmsans\fontsize{18pt}{18pt}\selectfont Nakiri Ayame với cách chữa trầm cảm tháng Năm\footnotemark} \\[8pt]
      % Dòng thứ tư: Nhân vật xuất hiện
       & {\caviar{\fontsize{14pt}{14pt}\selectfont Nhân vật: \emp{kuon1} \emp{ryuuhou1} \emp{achan} \emp{roboco2} \emp{ayame} \emp{subaru}}}
      \\[8pt]
      % Dòng cuối cùng: Thời gian diễn ra của tập (thường sẽ là ngày phát hành)
       & {\continuum\fontsize{16pt}{16pt}\selectfont Ngày phát hành: 02/05/2021}                               \\[18pt]
    \end{tabular}
  \end{adjustwidth}
\end{table}
\footnotetext{Đây không hẳn là một loại bệnh, mà đây là một hiện tượng tâm lý xảy ra với người mới đi làm hoặc sinh viên năm nhất khi họ cảm thấy mất động lực sau tháng đầu tiên đầy hứng khởi. Nguyên nhân chính của hội chứng này là sự kiệt quệ về tinh thần do căng thẳng, mất phương hướng hoặc không thích nghi được với môi trường làm việc/học tập, nên bất kỳ ai, kể cả những người đã làm việc lâu năm, đều có thể bị ảnh hưởng, đặc biệt là sau các giai đoạn làm việc căng thẳng hoặc thay đổi môi trường.}
%=========================================TRUYỆN CHÍNH========================================
\color{black}

\txtbx[2]
{{\color{vol} \uline{Tóm tắt:}} Một buổi sáng tháng Năm, Ayame - phát ra một lượng Nakirium cực mạnh - quyết tâm chữa ``bệnh trầm cảm tháng Năm'' cho mọi người bằng cách\dots\ quậy banh văn phòng. Từ việc ánh Subaru bay lên trần, cho tới làm Roboco quá nhiệt bốc khói và đập vỡ hai ô cửa sổ - thế mà cô vẫn bảo đó là ``vận động''. Sau màn hỗn loạn, cả nhóm bị A-chan phạt, còn Kuon và Miharu đánh giá thước phim là ``chưa đạt''. Cuối cùng, Roboco đề xuất một bài tập dưỡng sinh khá ổn, và Kuon cũng bị lôi vào tập cùng.}

Buổi sáng đầu tháng Năm mở đầu với không khí se lạnh còn sót lại của mùa xuân sắp hết, nhưng những tia nắng mặt trời đã bắt đầu ấm lên, báo hiệu mùa hè đang đến gần. Văn phòng của \hll\ hôm nay có vẻ náo nhiệt bất thường, không phải do công việc chất đống hay lịch trình dày đặc, mà vì một hiện tượng lạ vừa xuất hiện trên hệ thống giám sát của họ.

\begin{center}
  \uline{Phòng điều hành trung tâm.}
\end{center}

A-chan đứng trước màn hình hiển thị dữ liệu bất thường, mắt dán vào đường đồ thị vừa bật lên đột ngột, chạm ngưỡng cảnh báo đỏ. Cô chống cằm, lẩm bẩm với chính mình: ``Mức Nakirium tăng đột biến\dots\ có chuyện gì đang xảy ra sao?''

Vừa kịp suy nghĩ thêm, tiếng chân chạy vội vã vang lên từ hành lang. Một nhân viên kỹ thuật hớt hải chạy tới, cầm bảng dữ liệu trong tay, vừa thở dốc vừa báo cáo: ``Yuujin-san! Hệ thống vừa ghi nhận lượng Nakirium tăng vọt bất thường! Độ đậm đặc cao gấp 50 lần mức trung bình của cả tháng!''

``Cái gì? Chắc chỉ là lỗi hệ thống thôi\dots'' nàng Đeo kính nhướng mày, vẫn chưa thể tin được.

Nhưng cậu nhân viên kiên quyết lắc đầu, nhấn mạnh từng chi tiết: ``Không phải đâu ạ! Chúng tôi đã kiểm tra lại nhiều lần. Nguồn phát sóng mạnh nhất đến từ\dots\ văn phòng nơi Hikawa-san đang làm việc!''

Nghe đến đây, A-chan chỉ biết thở dài, lẩm bẩm bất lực: ``\hl{Lại là Nakiri-san nữa rồi\dots{} Cô ấy định bày trò gì nữa đây?}''

\begin{center}
  \uline{Văn phòng\\Cùng thời điểm.}
\end{center}

Tại văn phòng chính, bầu không khí cũng không kém phần sôi động. Kuon đang kiểm tra lại chiếc máy quay phim mà cậu đã dùng để quay Ayame từ nửa năm trước\footnote{Xem lại {\flaregothic ÉPISODE 81}, quyển số Sáu.}. Bên cạnh cậu là Miharu, bạn cùng lớp chuyên ngành lập trình, đứng nhìn màn hình với vẻ tò mò hứng thú. Và trên ghế sofa ở góc phòng, Ryuuhou đang buông điện thoại, tay chống cằm quan sát hai anh lớn với ánh mắt lấp lánh.

``Lại quay phim cho Ayame-san nữa hả, Onii-san?'' cô đặt điện thoại xuống, kéo ghế lại gần bàn làm việc. ``Lần này có gì đặc biệt hơn lần trước không?''

Kuon chẳng buồn ngước mắt khỏi máy quay, chỉ đáp gọn: ``Xem thì biết.''

Ngay lúc đó, Ayame - nhân vật chính của thước phim vừa rồi - bước vào với nụ cười rạng rỡ đặc trưng, tay vẫy chào hết mọi người: ``Lần này lại phiền anh quay giúp em thêm một thước phim nữa nhé!''

Miharu nghiêng đầu nhìn cậu Tổng, rồi quay sang Ayame: ``Vậy tại sao tôi cũng phải có mặt ở đây?''

``Thì anh phụ trách phần hậu kỳ mà. Kuon-senpai, máy quay đã sẵn sàng chưa ạ?'' nàng Quỷ sừng đứng ngay trước ống kính, hào hứng không khác gì những lần quay trước.

Kuon giơ máy lên, giơ ngón tay cái ra hiệu đã sẵn sàng. Ryuuhou ngồi bên cạnh mỉm cười, nhỏ đủ để chỉ Kuon nghe: ``Ayame-san lúc nào cũng tràn đầy năng lượng nhỉ. Mà kiểu năng lượng đó cũng lây nhiễm thật.''

Ngay trước khi Kuon bấm nút ghi hình, Ayame đột ngột quay sang Miharu với một yêu cầu khiến cậu bất ngờ: ``Anh có thể bế em lên một chút được không?''

``Ủa? Để làm gì?''

``Để anh ném em xuống sàn, nhẹ nhàng thôi, như em là một con thú nhồi bông ấy!'' cô đáp với đôi mắt sáng lên vẻ tinh nghịch.

Miharu há hốc miệng. Cậu lập trình viên chưa bao giờ nghĩ sẽ nhận được một đề nghị kỳ quặc như vậy từ một thành viên của \hll. Cậu chớp mắt vài lần, rồi nhìn sang Kuon như muốn tìm lời giải thích.

Nhưng cậu Tổng chỉ khoanh tay, gật đầu một cách bình thản: ``Đừng hỏi gì cả. Cứ làm đi. Cách em ấy bắt đầu đoạn phim năm ngoái cũng hao hao thế này.''

Ryuuhou ở bên cạnh bật cười thành tiếng, phá vỡ sự ngơ ngác của cậu bạn cùng lớp: ``Anh nói vậy làm em càng tò mò không chịu được!''

Cuối cùng, dù vẫn còn ngạc nhiên, Miharu cũng chỉ có thể làm theo lời Ayame. Cậu nâng cô lên một cách dễ dàng đến mức ngớ người - thật sự, Ayame nhẹ hều như một con thú bông. Cô cười khúc khích khi được nâng lên, càng làm cho bầu không khí trở nên vui vẻ.

``À, Kuon-senpai, anh có thể hướng máy quay xuống sàn được không?'' Ayame chỉ tay xuống đất.

Kuon thở dài, lẩm bẩm: ``Chả biết năm nay em ấy định làm trò gì nữa,'' rồi điều chỉnh ống kính hướng xuống sàn. Cậu trao đổi ánh mắt với Miharu, hai người cùng gật đầu.

``Ryuuhou, hô đi,'' cậu Tổng ra lệnh cho cô em gái, đáp lại cậu là một tiếng ``Vâng ạ\~{}'' đầy bùi tai của một con nhóc nghịch ngợm.

``Sẵn sàng\dots\ bắt đầu!''

\il

Cảnh quay mở đầu với góc nhìn hướng về bàn làm việc của A-chan, trông yên tĩnh đến lạ. Nhưng ai cũng biết sự bình yên này sẽ sớm bị phá vỡ.

Ayame, vẫn nằm trên tay Miharu như một con thú nhồi bông, giơ ngón tay ra hiệu. Và rồi\dots\ *BỤP!* cậu chàng ném nàng quỷ nhỏ xuống sàn không chút do dự. Cú tiếp đất phát ra tiếng động lớn, nhưng đáng ngạc nhiên là cô chẳng hề tỏ ra đau đớn. Ngược lại, máy quay ngay lập tức phóng to vào khuôn mặt tươi cười rạng rỡ của cô khi cô hét lên câu chào quen thuộc: ``\textbf{Yo-dazo! (Là tôi đây!)}''

\gif{ayame1}{1}{63}

Ryuuhou vỗ tay nhẹ nhàng từ phía sau máy quay, mỉm cười khen ngợi: ``\hl{Hay thật, Ayame-san nhập vai nhanh quá.}''

Nàng Oni không đứng dậy ngay mà bắt đầu xoay người trên sàn theo chiều ngược kim đồng hồ, trông giống hệt một con cá mắc cạn đang cố gắng bơi về lại nước. Cô tiếp tục màn diễn với giọng điệu nghịch ngợm: ``Mùa đông lạnh giá đã qua rồi, trời bây giờ đã bước vào xuân!''

Rồi bằng một động tác nhanh nhẹn, cô bật dậy và nhảy lên chiếc bàn gỗ gần đó, tiếp tục giả vờ là một con cá: đôi chân chổng lên trời, nằm sấp, hai cánh tay duỗi thẳng. ``\textbf{Mà đây còn là mùa cá hồi bơi ngược dòng nữa!}''

\gif{ayame2}{1}{78}

Từ sau máy quay, Kuon và Miharu đồng thanh thì thầm, giọng đầy bất lực: ``\hl{Mùa xuân đã đến từ Tết Nguyên Đán rồi em ơi\dots}''. Ngược lại, Ryuuhou cố nhịn cười, nghiêng người về phía sau trên ghế như thể vừa xem được đoạn hài hước nhất: ``\hl{Ôi trời ơi, Ayame-san làm em cười bụng đau luôn!}''

Cảnh chuyển tiếp sang Ayame nằm sấp trên ghế sofa, hai tay duỗi thẳng, lắc lư cả thân người. Cô ngẩng đầu lên, chớp mắt như vừa nhớ ra một sự thật quan trọng, rồi tuyên bố dõng dạc: ``Nhưng đây cũng là mùa mà con người bắt đầu có dấu hiệu lười biếng.''

\gif{ayame3}{1}{61}

``Đúng thật,'' Kuon và Miharu cùng gật gù đồng tình. Ryuuhou thì nhanh tay lấy điện thoại ra, quay lại màn hình máy quay như muốn lưu giữ khoảnh khắc đáng nhớ này.

Cảnh đột ngột cắt sang Ayame đứng bật dậy, cúi gập người, xoay hai tay theo vòng cung như đang bắt chước một loài vật: ``\textbf{Giống như những chú gấu ngủ đông dậy vậy!}''

\gif{ayame3}{62}{77}

Ngay lập tức, cảnh lại chuyển sang cô ngồi trên ghế xoay, quay một vòng đầy thần thái rồi chỉ thẳng vào ống kính: ``Vậy nên hôm nay, mình sẽ đi tìm những kẻ lười biếng đó và bắt họ phải vận động sức khỏe!''

\gif{ayame4}{1}{83}

Không dừng lại, cô nhảy lên chiếc bàn cạnh ghế sofa, lại tiếp tục giả vờ là con cá mắc cạn. ``Mình sẽ làm không chỉ họ, mà còn cho các bạn đang xem phải hét lên: `\textbf{Ôi quỷ thần ơi!}' vì cách thứ mà mình sẽ chỉ tới đây mà xem!''

\gif{ayame5}{1}{44}

Và đỉnh điểm của màn diễn đến khi cô lấy đà, hướng về phía cửa sổ với quyết tâm cao độ, hét lên khẩu hiệu cuối cùng: ``Nào, cùng mình xông ra phá vây nào!'' rồi lao thẳng vào cửa sổ, phá tung tấm kính ra ngoài, mà vẫn giữ nguyên động tác của con cá mắc cạn.

*XOẢNG!*

\gif{ayame5}{45}{72}

Từ bên trong văn phòng, Kuon và Miharu đổ mồ hôi hột sau pha hành động bất ngờ đó. Ryuuhou thì reo lên thích thú, chạy lại gần chỗ cửa kính vỡ: ``Ayame-san làm gì mà dữ dội thế? Cả ô cửa kính tan nát hết rồi kìa!''

Cậu bạn Lập trình quay sang Kuon, lẩm bẩm khó hiểu: ``Có nhất thiết phải làm đến mức đó không?''

Cậu Tổng thở dài, tặc lưỡi: ``Chịu em ấy luôn. Nhưng mà em ấy sẽ phải đền tiền thay cửa kính đấy.'' Cậu đưa máy quay cho Miharu.

``Sao vậy?''

``Giờ đến lượt ông cầm máy quay ghi hình tiếp đi.''

Miharu chẳng biết nói gì hơn, đành đồng ý. Ryuuhou đứng dậy, phủi nhẹ bụi trên quần, rồi đề xuất với ánh mắt hào hứng: ``Hay em cầm máy quay phụ cho? Em muốn học thêm về cách quay những cảnh thú vị thế này, nhìn nó phức tạp lắm.''

Kuon nhìn em gái, gật đầu đồng ý nhưng cũng nhắc nhở: ``Được thôi, nhưng đừng có mà làm trò gì lớn hơn Ojou đấy nhé. Em ấy mà dỗi vì bị cướp spotlight là anh không giúp đâu đấy.''

Cô em gái nheo mắt cười tinh nghịch, giơ tay lên như đang tuyên thệ: ``Em hứa!'' cô nhận lấy máy quay phụ từ tay Miharu, bắt đầu căn chỉnh ống kính với vẻ tập trung của một người đam mê nhiếp ảnh.

Miharu đứng một bên, chỉ biết lẩm bẩm với chính mình: ``\textit{Kiểu này chắc lại có đám người nhập viện vì quá liều Nakirium cho mà xem.}''

\ct

Ống kính hướng về Subaru đang rũ vai xuống bàn, ánh mắt lờ đờ như vừa cạn kiệt toàn bộ năng lượng. Cô thở một hơi dài như tiếng còi tàu, lẩm bẩm: ``Trời ạ, hôm nay chẳng muốn động tay động chân vào việc gì cả.''

Đối diện, Kuon cũng vươn vai uể oải, một tay bóp vai, tay kia xoa gáy như thể vừa trải qua ba đêm liên tục tăng ca. Cậu gật gù đồng tình: ``Công nhận. Tự dưng bước sang tháng Năm là thấy lười hẳn đi, không hiểu vì sao luôn.''

Ryuuhou thì đứng ở góc phòng, tay lướt điện thoại, cười nhẹ nhàng nhưng ẩn chứa sự tinh nghịch: ``Nghe hai anh chị nói thế, em cũng muốn dừng mọi thứ lại để nằm dài một lát thôi.'' Chiếc máy quay mà cô cầm hiện đang đặt ở phía trước chiếc TV.

Nhưng trước khi cả ba kịp hưởng thụ thêm chút sự lười biếng, Ayame đã nhô đầu ra từ phía sau, nụ cười trên môi toát lên vẻ của kẻ sắp thực hiện một kế hoạch phá hoại. Đôi mắt cô sáng rực, hai tay ôm chặt một cây chày nhôm lấp lánh ánh bạc.

\img{8-8a}

*BỤP BỤP*

Tiếng chân chạy vội vã vang lên khi nàng Quỷ nhỏ lao tới như một vận động viên bóng chày chuyên nghiệp, hét lớn: ``\textbf{Người đập bóng số một, Quỷ Oni đây!!}'' Rồi *VỤT!* một cú vung chày bất ngờ.

Cây chày nhôm quất mạnh vào lưng ghế của Subaru, không chỉ hất tung cô nàng Tomboy mà còn kéo đổ cả chiếc bàn bên cạnh. Máy quay rung mạnh bần bật, tạo cảm giác như vừa có cơn động đất nhỏ. Ryuuhou dù bất ngờ trước lực đánh của Ayame, vẫn giữ vững máy quay ổn định, khẽ nhận xét: ``Ayame-san đánh mạnh thật. Lực ấy mà trúng người thì chắc đi bệnh viện mất.''

Subaru - với tốc độ vượt xa trọng lực - bay vút lên cao, đầu va mạnh vào trần nhà kêu *CỐP!* một tiếng đau điếng, rồi rơi thẳng xuống ghế đối diện.

\img{8-8b}

Cùng lúc đó, Kuon nhanh chân nhảy lùi tránh được vụ hỗn loạn trong gang tấc. Cậu nhìn Ayame với vẻ mặt hoảng hốt: ``Ojou ơi, sao tự dưng lại hành hạ người ta như vậy?'' Cô nàng vác cây chày sau lưng, một tay chỉ thẳng vào hai người, mắt long lanh chẳng có chút hối lỗi mà còn hăng hái tuyên bố: ``Hai người này! Dừng ngay cái vẻ lười biếng kia lại!''

``Cái gì nữa vậy!?'' / ``Hả!?'' cả hai đồng thanh thảng thốt.

Nàng Quỷ sừng lùi vài bước, bẻ cổ kêu rắc rắc như một tay xã hội đen sắp đánh nhau, rồi chỉ thẳng vào hai người với dáng đứng đầy uy hiếp: ``Tôi đến đây để chữa căn bệnh trầm cảm tháng Năm của hai người!''

Kuon nghe vậy, mắt giật giật tức giận: ``Bọn anh có phải là lính mới đâu mà bị cái hội chứng đó!?'' Miharu cũng lên tiếng, nhướng mày lo ngại: ``Bọn anh cũng chẳng còn là sinh viên năm nhất nữa đâu\dots''

Nhưng Ayame bỏ ngoài tai mọi lời phàn nàn, cô rút từ sau lưng ra hai thanh kiếm gỗ, rồi cắn chặt cây chày vào miệng, đung đưa vung kiếm như một kiếm sĩ thực thụ. ``Lấy lại năng lượng nào!''

Cả Subaru và Kuon đồng thanh nổi cáu: ``\textbf{Cái này là quá tải cảm xúc chứ không phải chữa bệnh đâu!!}''

\gif{ayame6}{64}{122}

Nàng Vịt gầm lên, chộp lấy một hạt đậu khô trên bàn rồi ném mạnh về phía Ayame, như thể đang thực hiện phong tục đuổi quỷ của lễ hội Setsubun - vốn đã kết thúc từ ba tháng trước. ``\textbf{Nhận lấy này!!}''

Ayame hét toáng lên ngay khi thấy hạt đậu, mặt biến sắc vội vàng vứt cả chày lẫn kiếm, quay lưng bỏ chạy như thật sự bị trừ tà: ``\textbf{Eek! Đừng ném đậu vào tôi chứ! Phải chạy thôi!}'' Ryuuhou phá lên cười, vẫn giữ máy quay ổn định ghi lại khoảnh khắc: ``Ayame-san sợ hạt đậu còn hơn sợ ma nữa. Cảnh này chắc là khoảnh khắc vàng của thước phim này.''

Nàng Quỷ sừng lao nhanh ra cửa, để lại hai người vừa bị tấn công còn đứng hình chưa kịp hoàn hồn. Miharu đứng ngoài cửa quan sát, lẩm bẩm ngờ vực: ``\dots À, cái lễ Setsubun kết thúc từ lâu rồi đúng không?''

Kuon bỏ tay vào túi, liếc nhìn hạt đậu trên sàn rồi bình thản đáp: ``Tôi nghĩ ông nên hỏi em ấy kiếm hạt đậu đó ở đâu ra, chứ không phải hỏi lễ hết chưa. Mấy đứa ở đây có khái niệm thời gian quái đâu mà.''

Subaru liếc về phía cửa vừa biến mất của Ayame, thở dài ngao ngán: ``Thật không hiểu nổi cô ấy nghĩ gì.''

Tưởng mọi chuyện đã tạm lắng, sự yên tĩnh chẳng kéo dài được lâu. Nàng Oni lại ló đầu ra từ cửa, mặt cười toe toét như chưa hề có cuộc rượt đuổi vừa rồi: ``Bà thấy thích chứ?'' Cô bạn cùng Thế hệ đang vỗ đầu kiểm tra chấn động, nghe câu đó liền nổi điên hét lớn: ``\textbf{Im mồm đi!!}'' còn cậu Tổng thì chỉ khoanh tay, nhìn Ayame cười khổ trong khi lo lắng về những hậu quả tiếp theo.

\img{8-8c}

Cậu bạn Lập trình hỏi với giọng đều đều, như thể vừa chứng kiến quá nhiều chuyện vô lý trong một ngày: ``Cô ấy lúc nào cũng thế à, Kuon?'' Cậu Tổng thở dài, lắc đầu ngao ngán: ``Nhây nhây không bao giờ thay đổi. Chuyên phá văn phòng và bày những trò không thể tin nổi.''

Ryuuhou vừa tinh chỉnh ống kính máy quay, vừa xen vào: ``Em thấy Ayame-san có gì đó rất đáng yêu, dù có hơi điên một chút. Nhưng mà phá hai cửa sổ trong một buổi thì chi phí sửa chữa chắc không nhỏ đâu.''

Miharu vỗ vai Kuon, ánh mắt tỏ sự đồng cảm: ``Ngày nào cũng phải giải quyết những chuyện thế này chắc vất vả lắm nhỉ.'' Đáp lại, cậu Tổng chỉ thở dài. Nhưng trước khi kịp nói gì, cậu đột nhiên dừng lại, mắt nheo lại: ``Mà\dots\ ông có cảm thấy nóng trong phòng không?''

Miharu cũng nhướng mày: ``\dots Ừ, tự dưng thấy ấm hẳn lên. Nhiệt độ hôm nay mới có 15, 16 độ mà.'' Ryuuhou cũng đặt tay lên trán, khẽ nói: ``Ừ nhỉ, em cũng cảm thấy hơi nóng\dots\ có gì đó không ổn. Em không nhớ là hệ thống sưởi của công ty lại nóng tới mức thế--'' nhưng rồi chết điếng khi nhìn nghiêng sang trái..

Người anh trai đưa tay quạt quạt, rồi nhìn về phía nguồn nhiệt tỏa ra, khựng lại vài giây: ``Trời ơi, có chuyện rồi.''

Cả Miharu, Kuon, Ryuuhou và Subaru cùng dõi mắt về phía Roboco, người đang co giật từng đợt trên sàn ở góc phòng, cơ thể nóng lên như một chiếc CPU sắp quá tải. Hơi nóng bốc lên từ cô như ngọn lửa ngầm, nhuộm cả căn phòng màu cam rực rỡ.

\img{8-8d}

``\textbf{Lại có chuyện gì nữa đây!?}'' nàng Vịt thảng thốt, vội lùi ra sau vài bước.

Nàng Quỷ sừng đặt hai tay lên hông, gật gù đắc ý: ``Tôi đoán trước chuyện này sẽ xảy ra, nên đã làm Roboco-san quá nhiệt từ nãy giờ!'' Cô nói với giọng tự hào, không chút hối lỗi, khiến cả bốn người - cộng thêm một con robot đang bốc hỏa - đều đứng hình bất ngờ.

``Bà biết trước cái gì chứ!?'' / ``Có mà em đi phá người ta thì có!!'' Subaru và Kuon đồng thanh phản đối. Ryuuhou cũng lắc đầu cười khổ: ``Ayame-san, chị làm thế này thì ai mà chịu nổi nhiệt độ trong phòng chứ.''

Miharu đứng chống hông thở dài: ``Chả trách sao trong phòng nóng như lò nướng\dots\ Hóa ra là do cái con quỷ lắm tật này.'' Chưa kịp tra hỏi thêm, hệ thống cảnh báo trong phòng bật lên - hay đúng hơn, một cái nhiệt kế tự di chuyển lăn ra từ phòng thu, nói lớn: ``Cảnh báo: Nhiệt độ trong phòng đang tăng cao!'' rồi lẳng lặng biến mất từ cánh cửa.

Cùng lúc đó, thân thể nàng Người máy bùng lên một luồng lửa rực rỡ, khói trắng cuồn cuộn bốc lên từ bộ khung kim loại. Cô từ từ đặt tay lên ngực, giọng run run như vừa giác ngộ một chân lý thiêng liêng: ``Cả thân mình đang bùng cháy\dots\ Đây\dots\ có phải là tình yêu không?''

Ayame lập tức giơ hai tay tạo hình trái tim, cười toe toét. Nhưng ba người còn lại đồng loạt phản đối - Subaru quát ``\textbf{Không có đâu! Tất cả là lỗi của bà!}'', Kuon lắc đầu ``Rất tiếc là không đâu, Robo-PON à,'' còn Miharu chêm vào ``Xem quá nhiều phim ngôn tình rồi hả em?''

Ryuuhou vừa cười trừ vừa điều chỉnh tiêu cự máy quay: ``Roboco-san cũng bị lây cái chất ngôn tình từ Ayame-san rồi ạ? Em tưởng chỉ có mình chị ấy mới có năng khiếu đó chứ.''

\img{8-8e}

Căn phòng giờ thật sự biến thành phòng xông hơi miễn phí, nhiệt độ tăng nhanh chóng chóng mặt. Kuon hất cổ áo, thở hổn hển: ``Nóng quá! Không ổn rồi, phải mở cửa ra thôi!'' Nhưng trong khi mọi người hối hả tìm cách giải quyết, Ayame lại bò lên bàn, nằm sấp theo dáng một con cá hồi, hai tay vẫn giữ hình trái tim, miệng cười khoái chí: ``\textbf{Ha! Thành công rồi!!}''

``\textbf{Thành công cái gì chứ!?}'' cả ba người cùng thốt lên. Ryuuhou cũng theo giọng nói: ``Em thấy Ayame-san chỉ thành công trong việc biến văn phòng thành lò nướng thôi.''

Nàng Quỷ sừng chống cằm, nháy mắt đắc chí: ``Nóng thế này thì ai còn lười biếng được nữa, đúng không?''

\img{8-8f}

Nụ cười gian xảo của cô tỏa sáng giữa căn phòng rực lửa, nhưng chưa đầy ba giây sau, nàng Quỷ sừng đã lăn đùng ra sàn, thè lưỡi rên rỉ: ``Nóng\dots\ nóng quá\dots\ chết mất\dots''

Cậu Tổng khoanh tay, cúi xuống nhìn cô nàng đang giãy giụa như miếng thịt trên vỉ nướng: ``Ngu thật. Biết người mà không biết mình.'' Nàng Vịt chống hông, chẳng có chút thương cảm: ``Đúng là quỷ ngu.''

Nàng Người máy cũng cau mày: ``Gieo nhân nào gặt quả ấy thôi, em à.''

Ryuuhou tay máy quay ghi lại mọi thứ, tay kia quạt bản thân, giọng đùa giỡn nhưng vẫn giữ được sự bình tĩnh: ``Ayame-san mà cứ thế này thì sớm thành món sushi quỷ nướng trên bàn thôi.''

\img{8-8g}

Nhiệt độ vẫn tiếp tục tăng, mồ hôi lấm tấm trên trán Subaru. Cô quạt mạnh: ``Nếu thế này tiếp tục, tất cả chúng ta sẽ bị nướng chín hết!'' Nhưng Ayame bỗng nhổn dậy, trèo thẳng lên tường: ``Thử nói một câu chơi chữ dở tệ vào là nhiệt độ sẽ hạ xuống ngay! Tôi đảm bảo!''

Kuon nhìn cô với ánh mắt hoài nghi: ``Nói bừa vừa thôi em ơi.''

Nàng Quỷ sừng giơ ngón tay cái lên trời, cười tự mãn: ``Cứ tin tôi đi! Chắc chắn hiệu nghiệm!''

Subaru và Miharu nhìn nhau ngờ vực, cuối cùng Subaru hít một hơi thật sâu, rồi nói: ``Muối này mặn mòi, muốn nói một lời?'' Ngay lập tức, Ayame lăn lộn trên sàn cười ngất ngây, như vừa nghe được câu chuyện hài hước nhất thế giới. Subaru há hốc miệng: ``Cô ấy thật sự thấy câu đó buồn cười sao!?''

Miharu chép miệng: ``Nhạt thế mà cũng cười được, tài thật.'' Kuon nhún vai: ``Thì cô ấy lúc nào cũng dễ cười mà. Nhưng mà cái trò này thì giải quyết được gì đâu.'' Ryuuhou vừa quay phim vừa cười theo: ``Em dám cá Ayame-san mà làm diễn viên hài thì chỉ có mình cô ấy cười thôi, khán giả thì chẳng ai hiểu nổi điểm hài ở đâu.''

Ngay lúc đó, Roboco bỗng run lên bần bật, giọng lắp bắp: ``C-Câu chuyện\dots\ lạnh quá\dots'' một luồng hơi lạnh đột ngột tỏa ra từ cô, làm nhiệt độ phòng hạ xuống nhanh chóng.

\img{8-8i}

Miharu và Kuon nhìn nhau với vẻ không thể tin nổi: ``Thế cũng được thật sao!?'' trong khi Subaru lau mồ hôi trên trán, thở phào nhẹ nhõm: ``Mà dù sao thì chúng ta cũng thoát chết rồi.''

Ayame nhảy cẫng lên, chạy đến chỗ Subaru, đôi mắt lấp lánh: ``Bà thấy cách của tôi hay không?'' nhưng không để ai nêu cảm nhận, Roboco đã vẫy tay gọi cô: ``Lại đây nào, Ayame-chan\~{}'' với nụ cười trên môi chẳng mấy thiện cảm.

Kuon lập tức nói với Ryuuhou: ``Hướng máy về phía Ayame đi, sắp có chuyện hay rồi.'' Lập tức, Ryuuhou cũng nhanh trí điều chỉnh góc quay, giọng hào hứng mà vẫn tỉnh táo: ``Phải quay cảnh này ở chế độ chậm mới thật sự ấn tượng.''

Nàng Quỷ sừng chạy đến chỗ tiền bối Người máy, tò mò hỏi: ``Sao vậy chị?''

Roboco không trả lời, chỉ quay sang cậu Tổng: ``Đưa cho em cái chày đây.''

Kuon đáp không do dự: ``Rất sẵn lòng,'' rồi đặt cây chày gỗ ngay vào tay Roboco, cùng Miharu chạy tót về phía cô em gái của mình.

Một giây im lặng trôi qua.

Rồi\dots *BỘP!!!*

Ayame bay khỏi mặt đất, văng thẳng về phía cửa sổ với tốc độ kinh ngạc, vỡ tung thêm một ô kính nữa, biến mất thành một bóng đen nhỏ ngoài trời xanh.

``\textbf{AAAAAAAAAAAAAA\~{}!!!}''

Ryuuhou há hốc miệng khi theo dõi qua ống kính, khẽ thốt lên: ``Chà\dots\ Roboco-san đánh bóng chuyên nghiệp thật. Lực ấy mà vào thi đấu chắc vô địch.'' Miharu và Roboco nhìn nhau, rồi cùng thốt lên đầy phấn khởi: ``\textbf{Một cú đánh hoàn hảo! Quỷ sừng bị loại khỏi sân!}''

Kuon khoanh tay, nhìn đống kính vỡ dưới sàn, thở dài: ``Lại thêm một ô cửa sổ nữa phải sửa\dots\ Nhưng mà\dots'' Cậu và Subaru quay sang nhau, đồng thanh chống hông, giọng chất chứa sự thỏa mãn: ``Cho đáng đời đi.''

Ryuuhou bấm nút dừng ghi hình, giơ máy quay lên với vẻ hài lòng: ``Cảnh quay này nếu đăng lên mạng chắc chắn sẽ thành viral, đắt hàng làm meme lắm.'' Cậu bạn Lập trình thì chỉ biết đứng yên nhìn mọi thứ, lắc đầu cười trừ trước những diễn biến khó tin vừa xảy ra.

\begin{center}
  \uline{Đoạn phim kết thúc.}
\end{center}

%==========================================TRUYỆN PHỤ=========================================
\omk

\begin{center}
  \uline{Sau buổi ghi hình.}
\end{center}

Kẻ đứng sau tất cả vụ hỗn loạn - Ayame - vẫn đang tràn đầy Nakirium, nhưng lần này không phải vì phấn khích, mà vì vừa bị A-chan mắng thậm tệ sau khi phá tan vô số cửa kính văn phòng. Dù bị chửi cho nhũn cả mặt, cô nàng vẫn giữ được vẻ mặt thản nhiên như thể chuyện đó chẳng liên quan gì đến mình, thỉnh thoảng còn vén nhẹ tóc mái một cách tự nhiên, chẳng khác nào một ngôi sao đang bình thản nhận phỏng vấn.

Roboco, người đã vô tình làm thêm một cái cửa sổ nữa vỡ tung sau cú đánh Ayame bay ra cửa, cũng phải nhận lỗi chung. Kết quả, cả hai đều bị ghi tên vào sổ phạt lương để đền tiền sửa chữa - một cảnh tượng quá quen thuộc đến nỗi A-chan chỉ cần nhắc một câu, nhân viên tài chính đã tự động cập nhật số tiền phạt vào hệ thống, chẳng cần hỏi thêm gì. 

Trong khi đó, Ryuuhou đang ngồi xếp các mảnh kính vỡ vào thùng rác, tay làm việc đều đặn nhưng miệng vẫn trêu chọc Subaru để giải tỏa bầu không khí căng thẳng: ``Chị đừng quét mạnh quá, Subaru-san, kẻo lát lại đập vỡ thêm cái gì nữa thì lại phải thêm cửa sổ đền đấy.'' Cô bé nhấc một mảnh kính to lên, vứt nhẹ vào thùng, giọng vẫn cười tinh nghịch, dù hành động của mình khá cẩn thận. Subaru nổi điên, quét mạnh mấy nhát chổi vào chân Ryuuhou để đùa: ``Còn không phải do hai người kia gây ra à! Mà chị còn bị bay lên trần nhà nữa cơ, sao không có ai đền cho chị chứ?''

Ở góc phòng đối diện, Kuon và Miharu cùng ngồi xem lại thước phim vừa quay xong trên màn hình laptop. Cậu bạn Lập trình ngả người ra ghế, khoanh tay bình phẩm: ``Cái này trông lộn xộn quá.''

Cậu Tổng quay sang nhìn cậu bạn, mắt hơi nheo lại: ``Lộn xộn sao? Tôi thấy quay cũng ổn mà?''

Miharu tua lại đoạn đầu, chỉ vào màn hình: ``Ban đầu chúng ta thống nhất là quay đoạn phim về cách chữa bệnh trầm cảm tháng Năm, đúng không? Rồi từ đầu đến cuối, có thấy cách chữa đâu đâu? Toàn là nhảy nhót, đánh người, rồi phá công ty tan hoang.''

Kuon thở dài, chống cằm xem lại đoạn Ayame lao qua cửa sổ: ``Thực ra em ấy cũng cố gắng làm mọi người vận động nhiều thật, nhưng cách làm thì\dots\ hơi khác một chút.''

``Không phải hơi khác, là chệch hướng hoàn toàn luôn,'' Miharu cười khổ. ``Ban đầu còn đi đúng đường, sau đó càng lúc càng lệch, cuối cùng vỡ mất hai cái cửa sổ. Báo hại công ty tổn thất không ít.''

``\dots Nhìn kỹ lại thì cậu nói đúng. Nếu thay đổi cách thực hiện thì chắc chắn sẽ hợp lý hơn nhiều.''

``Đấy là tôi chỉ nhận xét vậy thôi,'' Miharu ngước mắt khỏi màn hình, nhìn sang Kuon, ``còn ông thì thấy sao? Phải làm lại không?''

Kuon ngả người ra sau, tặc lưỡi một cái rồi kết luận dứt khoát: ``Chưa đạt tiêu chuẩn. Phải quay lại.'' Cậu chẳng nói gì thêm, chỉ gật đầu đồng tình. Cả hai ngồi chờ mọi người dọn dẹp xong, chuẩn bị cho cuộc họp đánh giá chung để cùng sửa lỗi cho thước phim.

\ct

Sau khi tất cả tập trung lại ở phòng nghỉ, Kuon mở đầu cuộc họp bằng những nhận xét về lỗ hổng của thước phim mà cậu cùng Ayame và Miharu thực hiện. Nhân vật chính của thước phim chớp chớp mắt ngơ ngác sau khi cậu Tổng đánh giá: ``Như vậy là\dots\ đoạn phim chưa ổn ạ?''

Kuon lắc đầu, tỏ rõ sự không hài lòng: ``Chưa. Quá lộn xộn, và đã đi sai hoàn toàn so với mục tiêu ban đầu mà chúng ta đặt ra.''

Subaru nghe vậy liền nhảy vào bổ sung, mặt vẫn còn cau có: ``Đúng đó! Còn phá gần như tanh bành cả văn phòng nữa! Suýt nữa là em thật sự phải đi bệnh viện kiểm tra chấn thương sọ não rồi. Giờ mà bảo cầm chày đập vào người là để vận động, ai mà tin được chứ!?''

Ryuuhou thì chỉ nở một nụ cười gượng sau khi đích thân cô đã ghi lại toàn bộ sự việc lúc đó: ``Em phải đồng ý với Onii-san. Có nhiều cách để có thể khuyến khích người ta vận động được mà. Cách của chị thực sự rất nguy hiểm đấy ạ.''

Thấy mọi người đang bàn tán sôi nổi, Roboco ngồi ở ghế xa - dường như không hiểu chuyện gì xảy ra - liền giơ tay: ``Mọi người đang nói về chuyện gì thế? Tóm tắt lại cho chị được không?'' Cô em gái của cậu Tổng đáp lại: ``Ủa, Roboco-san không biết ạ?''

``Không,'' nàng Người máy lắc đầu, giọng có chút hơi bất mãn. ``Chị chỉ biết Ayame-chan bảo là `cho em mượn chị có chút việc', và thế là cứ thế tự ý tăng nhiệt độ của chị thôi.''

Ayame, Kuon và Miharu lần lượt kể lại toàn bộ câu chuyện: từ lúc họ lên ý tưởng quay một đoạn phim ngắn về căn bệnh trầm cảm tháng Năm, cho đến những tình tiết bất ngờ mà Ayame tự ý thêm vào khiến câu chuyện đi hoàn toàn sai hướng.

Roboco nghe xong mới gật gù hiểu ra: ``À\dots\ ra là hội chứng mệt mỏi tháng Năm á? Cũng đến mùa này rồi nhỉ?''

Ryuuhou vừa dọn xong mảnh kính cuối cùng, ngồi xuống ghế cạnh cô nàng Rô-bốet rồi nhún vai nhận xét: ``Em thấy hội chứng đó thường chỉ có mấy bạn mới đi làm hay bị thôi. Cả đám mình ở đây làm cũng gần hai năm rồi, tinh thần vững như núi, có mà lo gì.''

Kuon khoanh tay đồng tình: ``Đấy, tự dưng bày ra vụ này hơi vô lý quá. Có khi em nên đi tìm A-chan - người đúng là đang bị căng thẳng thật - để làm ví dụ thì còn hợp lý.'' Miharu cũng nhún vai, đơn giản hóa vấn đề: ``Cách chữa cái đó cũng đơn giản mà, chỉ cần ra ngoài vận động thường xuyên, đi dạo quanh công viên hay chơi thể thao gì đó là khỏe lại ngay thôi.''

Nàng Quỷ sừng lập tức phản đối, chỉ vào bản thân: ``Thì em cũng đang thúc giục mọi người vận động mà! Còn không phải vừa mới bắt Subaru-chan phải chạy nhảy đó sao!''

Cậu Tổng liền nheo mắt: ``Nhưng ai lại cầm chày gỗ mà thúc người ta vận động bao giờ? Làm vậy không phải thúc vận động, là cố tình gây thương tích đó!''

Giữa lúc cả nhóm đang tranh luận về cách làm sai lầm của Ayame, Roboco ngồi ngoài bỗng reo lên, mặt rạng rỡ: ``Em có một cách hay đây! Bảo đảm mọi người đều được vận động đúng cách!'' Cả nhóm lập tức quay đầu nhìn về phía cô, tò mò không biết cô định bày trò gì nữa.

Ryuuhou - vốn từng dẫn dắt ban nhạc underground của mình qua bao buổi tập luyện và trình diễn - bỗng khẽ nhếch mép, nói với giọng điệu của một người đã quen thuộc với việc điều phối nhóm: ``Roboco-san có ý tưởng gì, thử chia sẻ đi. Em có kinh nghiệm dẫn dắt nhóm tập luyện, nếu cần em sẽ giúp giữ nhịp cho mọi người.''

Kuon thở dài, khoanh tay chờ đợi: ``Được thôi, để xem ý tưởng mới của em có điên rồ hơn Ayame không\dots''

\ct

``\dots Cách này\dots\ cũng được thật à?'' Kuon nhăn mặt nhìn cảnh tượng trước mắt.

Miharu cũng nghiêng đầu, vừa tò mò vừa khó hiểu: ``Tôi cũng chưa bao giờ nghe về cách tập thể dục này. Trông nó\dots\ kỳ lạ thật.''

Trước mắt hai cậu con trai là bốn cô gái đang vung tay múa chân hết sức nhiệt tình. Hai tay giơ cao, thân thể lắc lư qua lại theo một nhịp điệu kỳ lạ, trông giống một bài dưỡng sinh nhưng có phần hơi quá khích. Subaru thì loạng choạng cố bắt kịp nhịp, Ayame hăng hái đến mức như sắp bay khỏi mặt đất, Roboco điều khiển chuyển động một cách máy móc, còn Ryuuhou đang đứng giữa, giọng hô nhịp đều đều nhưng vẫn giữ được nụ cười nghịch ngợm. Cô bé nhanh chóng nắm bắt động tác, thậm chí còn dẫn dắt cả nhóm đi đúng nhịp: ``Một, hai, ba, bốn! Lên nào các chị, đừng chùn bước, xong buổi này em dẫn đi uống trà sữa nhé!'' cô vừa nói vừa khéo léo điều chỉnh tốc độ cho Subaru chậm lại một chút, tạo thêm cảm giác gần gũi.

\gif{ayame7}{1}{146}

``Nhưng mà trông họ giống mấy con giun bị giày xéo quá nhỉ\dots'' Kuon đổ mồ hôi hột, đứng ngoài quan sát. Miharu nhún vai, cũng cười: ``Thôi thì cũng được, ít nhất là mọi người đều đang vận động thật, đúng với mục tiêu ban đầu mà.''

Nàng Oni vẫn giữ nguyên nụ cười rạng rỡ, vừa tập vừa cười khanh khách: ``Bài tập này hay quá! Sao em không biết sớm hơn, nếu biết thì đã rủ mọi người tập từ lâu rồi!''

Nàng Vịt thở hổn hển, tay run rẩy hạ xuống: ``Chúng ta nghỉ được chưa? Tôi chóng mặt quá rồi, sắp ngất đi thật\dots''

Nàng Người máy lập tức khoanh tay, đáp kiên quyết: ``Còn phải tập thêm mười phút nữa! Vận động nhiều vào mới khỏe người, chứ vừa mới tập được bao lâu mà đòi nghỉ!''

Nàng Trưởng nhóm cũng gật đầu phụ họa, dù miệng cười nhưng động tác vẫn đều đặn: ``Subaru-san cố lên! Sắp xong rồi, tập hết thì em dẫn mọi người đi ăn tráng miệng ở cửa hàng bánh ngọt gần đây nhé!'' cô vừa hô nhịp vừa khéo léo thúc đẩy mọi người, thể hiện sự trưởng thành hơn tuổi của mình.

Hai người con trai chỉ đứng quan sát, lâu lâu lại chụp vài tấm hình lưu lại khoảnh khắc độc nhất vô nhị này. Cả hai đều không thể nín cười, nhất là khi họ phải để một cô bé mười lăm tuổi dẫn dắt cả nhóm.

Bất chợt, Ayame lóe lên ý tưởng, quay ngoắt lại chỉ tay về phía hai người đàn ông: ``A! Hay là hai anh cũng vào tập cùng bọn em đi! Cùng vận động cho khỏe!''

Miharu ngay lập tức xua tay từ chối: ``À thôi, tí nữa anh có hẹn bạn tập bóng rổ rồi, các em tập tiếp đi nhé.'' Cô nàng lập tức chuyển mục tiêu sang Kuon: ``Vậy thì Kuon-senpai vào cùng bọn em! Anh cũng cần vận động mà!''

Ryuuhou nhìn anh trai với vẻ tinh nghịch: ``Onii-san, em mà còn tập được thì anh không lý nào chùn bước đâu nhỉ? Làm gương cho em đi chứ!'' cô nói thêm, như một lời thách thức nhẹ nhàng. Kuon ngập ngừng, liếc nhìn cảnh tượng kỳ lạ rồi thở dài: ``Ờm\dots\ được thôi.''

Dù ban đầu có hơi miễn cưỡng, cậu bước vào vòng tròn, giơ tay lên theo nhịp độ của mọi người. Lúc đầu cậu thấy kỳ quặc, nhưng dần dần khi hòa vào được nhịp điệu, cậu lại thấy cũng không tệ.

``\textit{Dù gì thì, ít nhất mọi người cũng đang vận động thật. Cũng coi như đạt được mục tiêu chữa bệnh trầm cảm tháng Năm rồi nhỉ\dots}'' cậu nghĩ bụng, ``\textit{\dots có hơi kỳ lạ một chút, nhưng cũng đáng thôi.}''

\end{document}